% GENERATED FILE. Edit canonical lessons, then run npm run generate:books.
% canonical-source-hash: fnv1a64:df1396007551dca8
% canonical-lessons: GE-C392-geraet, GE-C392-akku, GE-C392-aufladen, GE-C392-ladegeraet, GE-C392-stecker

\chapter{Geräte laden}
\label{ch:ge-gerate-laden}
\hlchaptermodality{392}

\begin{chapteropening}
\textbf{By the end of this chapter:} \emph{I can talk about devices and appliances, batteries, charging and plugs.}

Complete the last lesson of chapter 392: I can talk about devices and appliances, batteries, charging and plugs.
\end{chapteropening}

\section[das Gerät]{das Gerät — the device, the appliance}

\label{lesson:GE-C392-geraet}



\begin{quote}
Before the new one: say the German for to acquire, then the German for to insure.
\end{quote}

\subsection*{You'll want to know: das Gerät}
\textbf{das Gerät} — \textquotedblleft{}the device, the appliance\textquotedblright{}.

\begin{quote}
In der Küche stehen viele elektrische Geräte.
\end{quote}

\begin{quote}
\emph{There are many electrical appliances in the kitchen.}
\end{quote}

Plural: \emph{die Geräte}. In a gym, \emph{die Geräte} are the machines.

\begin{cousinweb}
\emph{Gerät} is inherited from Middle High German \emph{geræte} \textquotedblleft{}consultation; equipment\textquotedblright{} and Old High German \emph{girāti} \textquotedblleft{}advice\textquotedblright{} (EtymDB), of the family of \emph{Rat} \textquotedblleft{}advice\textquotedblright{}: what you need to provide for a task.
\end{cousinweb}

\begin{grammarlens}[title={What you've built}]
One more word for telling, explaining and describing at B1.
\end{grammarlens}

\subsection*{Your turn}
\begin{itemize}
\raggedright
  \item \emph{Say it:} \emph{das Gerät}
  \item \emph{Say it:} \emph{das Gerät}, once more
  \item \emph{Say it:} a sentence of your own with \emph{das Gerät}
  \item \emph{Recall:} say the German for to acquire, then the German for to insure, then say \emph{das Gerät} again
\end{itemize}

\begin{tcolorbox}[breakable,colback=teal!4,colframe=teal!35!black,title={Before you move on}]
What does das Gerät mean? (\textquotedblleft{}The device, the appliance\textquotedblright{}.) Say it once more.
\end{tcolorbox}
\section[der Akku]{der Akku — the rechargeable battery}

\label{lesson:GE-C392-akku}



\begin{quote}
Before the new one: say the German for to insure, then the German for the device.
\end{quote}

\subsection*{You'll want to know: der Akku}
\textbf{der Akku} — \textquotedblleft{}the rechargeable battery\textquotedblright{}.

\begin{quote}
Mein Akku ist fast leer.
\end{quote}

\begin{quote}
\emph{My battery is almost flat.}
\end{quote}

Plural: \emph{die Akkus}. Short for \emph{der Akkumulator}.

\begin{cousinweb}
\emph{Akku} is short for \emph{Akkumulator}, a learned formation on Latin \emph{accumulāre} \textquotedblleft{}to heap up, to accumulate\textquotedblright{} (Lewis \& Short), built on \emph{cumulus} \textquotedblleft{}a heap\textquotedblright{}: it stores up electricity.
\end{cousinweb}

\begin{grammarlens}[title={What you've built}]
One more word for telling, explaining and describing at B1.
\end{grammarlens}

\subsection*{Your turn}
\begin{itemize}
\raggedright
  \item \emph{Say it:} \emph{der Akku}
  \item \emph{Say it:} \emph{der Akku}, once more
  \item \emph{Say it:} a sentence of your own with \emph{der Akku}
  \item \emph{Recall:} say the German for to insure, then the German for the device, then say \emph{der Akku} again
\end{itemize}

\begin{tcolorbox}[breakable,colback=teal!4,colframe=teal!35!black,title={Before you move on}]
What does der Akku mean? (\textquotedblleft{}The rechargeable battery\textquotedblright{}.) Say it once more.
\end{tcolorbox}
\section[aufladen]{aufladen — to charge (a battery); to load}

\label{lesson:GE-C392-aufladen}



\begin{quote}
Before the new one: say the German for the device, then the German for the rechargeable battery.
\end{quote}

\subsection*{You'll want to know: aufladen}
\textbf{aufladen} — \textquotedblleft{}to charge (a battery); to load\textquotedblright{}.

\begin{quote}
Ich lade mein Handy jede Nacht auf.
\end{quote}

\begin{quote}
\emph{I charge my mobile every night.}
\end{quote}

Separable and strong: \emph{er lädt auf}, \emph{er hat aufgeladen}. Also to load: \emph{das Gepäck aufladen}.

\begin{cousinweb}
\emph{aufladen} is built from \emph{auf} \textquotedblleft{}up, on\textquotedblright{} and \emph{laden} \textquotedblleft{}to load\textquotedblright{}, inherited from Old High German \emph{hladan} and Proto-Germanic \emph{hlaþaną} (EtymDB): you load the battery up.
\end{cousinweb}

\begin{grammarlens}[title={What you've built}]
One more word for telling, explaining and describing at B1.
\end{grammarlens}

\subsection*{Your turn}
\begin{itemize}
\raggedright
  \item \emph{Say it:} \emph{aufladen}
  \item \emph{Say it:} \emph{aufladen}, once more
  \item \emph{Say it:} a sentence of your own with \emph{aufladen}
  \item \emph{Recall:} say the German for the device, then the German for the rechargeable battery, then say \emph{aufladen} again
\end{itemize}

\begin{tcolorbox}[breakable,colback=teal!4,colframe=teal!35!black,title={Before you move on}]
What does aufladen mean? (\textquotedblleft{}To charge (a battery); to load\textquotedblright{}.) Say it once more.
\end{tcolorbox}
\section[das Ladegerät]{das Ladegerät — the charger}

\label{lesson:GE-C392-ladegeraet}



\begin{quote}
Before the new one: say the German for the rechargeable battery, then the German for to charge.
\end{quote}

\subsection*{You'll want to know: das Ladegerät}
\textbf{das Ladegerät} — \textquotedblleft{}the charger\textquotedblright{}.

\begin{quote}
Ich habe mein Ladegerät im Hotel vergessen.
\end{quote}

\begin{quote}
\emph{I left my charger at the hotel.}
\end{quote}

Plural: \emph{die Ladegeräte}. The cable alone is \emph{das Ladekabel}.

\begin{cousinweb}
\emph{Ladegerät} joins \emph{laden} \textquotedblleft{}to load\textquotedblright{}, Old High German \emph{hladan} (EtymDB), to \emph{Gerät} \textquotedblleft{}device\textquotedblright{}, Middle High German \emph{geræte} \textquotedblleft{}equipment\textquotedblright{} (EtymDB).
\end{cousinweb}

\begin{grammarlens}[title={What you've built}]
One more word for telling, explaining and describing at B1.
\end{grammarlens}

\subsection*{Your turn}
\begin{itemize}
\raggedright
  \item \emph{Say it:} \emph{das Ladegerät}
  \item \emph{Say it:} \emph{das Ladegerät}, once more
  \item \emph{Say it:} a sentence of your own with \emph{das Ladegerät}
  \item \emph{Recall:} say the German for the rechargeable battery, then the German for to charge, then say \emph{das Ladegerät} again
\end{itemize}

\begin{tcolorbox}[breakable,colback=teal!4,colframe=teal!35!black,title={Before you move on}]
What does das Ladegerät mean? (\textquotedblleft{}The charger\textquotedblright{}.) Say it once more.
\end{tcolorbox}
\section[der Stecker]{der Stecker — the plug}

\label{lesson:GE-C392-stecker}



\begin{quote}
Before the new one: say the German for to charge, then the German for the charger.
\end{quote}

\subsection*{You'll want to know: der Stecker}
\textbf{der Stecker} — \textquotedblleft{}the plug\textquotedblright{}.

\begin{quote}
Im Ausland passt mein Stecker oft nicht.
\end{quote}

\begin{quote}
\emph{Abroad my plug often does not fit.}
\end{quote}

Plural: \emph{die Stecker}. \emph{Den Stecker ziehen} — to pull the plug.

\begin{cousinweb}
\emph{Stecker} is formed with \emph{-er} on \emph{stecken} \textquotedblleft{}to put in, to stick\textquotedblright{}, inherited from Old High German \emph{steckōn} (EtymDB): the thing you stick in.
\end{cousinweb}

\begin{grammarlens}[title={What you've built}]
One more word for telling, explaining and describing at B1.
\end{grammarlens}

\subsection*{Your turn}
\begin{itemize}
\raggedright
  \item \emph{Say it:} \emph{der Stecker}
  \item \emph{Say it:} \emph{der Stecker}, once more
  \item \emph{Say it:} a sentence of your own with \emph{der Stecker}
  \item \emph{Recall:} say the German for to charge, then the German for the charger, then say \emph{der Stecker} again
\end{itemize}

\begin{tcolorbox}[breakable,colback=teal!4,colframe=teal!35!black,title={Before you move on}]
What does der Stecker mean? (\textquotedblleft{}The plug\textquotedblright{}.) Say it once more.
\end{tcolorbox}
